% ============================================================
% InternNova Week 3 - Tasks 3 & 4
% Malware, Phishing & Security Awareness
% No cover page: prepared to continue from Tasks 1 & 2.
% Pattern follows the Week 2 report structure:
% Task -> topic -> explanation -> diagram -> cybersecurity
% connection -> security risk/protection -> reference basis ->
% reference/standard -> learning video -> practical/check.
% ============================================================

\documentclass[12pt,a4paper]{article}

% ---------- Packages ----------
\usepackage[margin=1in]{geometry}
\usepackage{setspace}
\usepackage{graphicx}
\usepackage{float}
\usepackage{booktabs}
\usepackage{array}
\usepackage{xcolor}
\usepackage{hyperref}
\usepackage{placeins}
\usepackage{needspace}
\usepackage{caption}
\usepackage{fancyhdr}
\usepackage{enumitem}
\usepackage{titlesec}
\usepackage{longtable}

% ---------- Formatting ----------
\onehalfspacing
\hypersetup{colorlinks=true,linkcolor=black,urlcolor=black,citecolor=black}
\setlength{\emergencystretch}{3em}
\sloppy
\urlstyle{same}
\Urlmuskip=0mu plus 1mu

\pagestyle{fancy}
\setlength{\headheight}{15pt}
\fancyhf{}
\lhead{InternNova -- Week 3}
\rhead{Tasks 3 \& 4: Phishing \& Malware Prevention}
\cfoot{\thepage}

\titleformat{\section}
  {\Large\bfseries}{\thesection}{0.7em}{}
\titleformat{\subsection}
  {\large\bfseries}{\thesubsection}{0.7em}{}
\titleformat{\subsubsection}
  {\normalsize\bfseries}{\thesubsubsection}{0.7em}{}

\setlist[itemize]{itemsep=3pt,topsep=4pt}
\setlist[enumerate]{itemsep=3pt,topsep=4pt}

% ---------- Reusable reference blocks ----------
\newcommand{\referenceblock}[2]{%
\begin{samepage}
\noindent\textbf{Reference basis:} #1\par
\smallskip
\noindent\url{#2}\par
\end{samepage}
\par\vspace{0.25cm}
}

\newcommand{\paperblock}[2]{%
\begin{samepage}
\noindent\textbf{Reference / Standard:} #1\par
\smallskip
\noindent\url{#2}\par
\end{samepage}
\par\vspace{0.25cm}
}

\newcommand{\videoblock}[3]{%
\begin{samepage}
\noindent\textbf{Learning Video -- S.No. #1:} #2\par
\smallskip
\noindent\url{#3}\par
\end{samepage}
\par\vspace{0.25cm}
}

\begin{document}
\setcounter{page}{1}

% ============================================================
% TASK 3
% ============================================================

\section{Task 3: Identify a Phishing Message}

\subsection{Introduction}

Phishing messages are designed to appear trustworthy while attempting
to persuade a recipient to click a link, open an attachment, provide
sensitive information, or perform another action. A safe analysis
should therefore focus on observable warning signs rather than on
interacting with the suspicious message.

For this task, two publicly documented phishing examples from
security-awareness resources were reviewed. The examples are used only
for educational analysis. Suspicious links, attachments, phone numbers,
and websites should not be opened or contacted during the analysis.

\subsection{Phishing Message Analysis Process}

A phishing message can be examined systematically by checking the
sender, links, language, requested information, and other unusual
content. After identifying the warning signs, the safest response is
to verify independently, report the message when appropriate, and
delete it.

\begin{figure}[H]
\centering
\includegraphics[width=0.48\textwidth]{images/phishing_analysis_flow.png}
\caption{Phishing Message Analysis Process}
\label{fig:phishing-analysis}
\end{figure}
\FloatBarrier

\subsection{Phishing Identification Checklist}

The following checklist groups common indicators that can be reviewed
without interacting with a suspicious message.

\begin{figure}[H]
\centering
\includegraphics[width=0.98\textwidth]{images/phishing_identification_checklist.png}
\caption{Phishing Identification Checklist}
\label{fig:phishing-checklist}
\end{figure}
\FloatBarrier

\subsection{Public Example 1: Fake MetaMask Phishing Email}

The U.S. Federal Trade Commission (FTC) published examples of fake
MetaMask and PayPal messages. The MetaMask message falsely claimed that
the recipient's cryptocurrency wallet was blocked and created pressure
to act quickly by clicking a link and updating the wallet. The FTC
warned that unexpected messages asking users to act quickly, click a
link, or call a number are common phishing patterns.

\begin{table}[H]
\centering
\small
\renewcommand{\arraystretch}{1.35}
\begin{tabular}{|p{4.0cm}|p{9.2cm}|}
\hline
\textbf{Indicator} & \textbf{Analysis}\\
\hline
Suspicious sender & The message impersonates a trusted service. The sender identity should not be accepted only because the message uses a familiar company name.\\
\hline
Suspicious link & The message asks the recipient to click a link to update or resolve the wallet problem. The link should not be opened.\\
\hline
Urgent / threatening language & It claims the wallet is blocked and suggests that the recipient must act quickly or risk losing cryptocurrency.\\
\hline
Spelling / grammar & Poor writing is not required for a message to be phishing. Modern phishing messages may be professionally written, so grammar alone is not a reliable test.\\
\hline
Request for personal information & The requested wallet update is intended to make the recipient interact with a fraudulent destination and potentially expose sensitive information.\\
\hline
Other suspicious element & The problem is unexpected and creates pressure to solve it through the message itself instead of independently verifying the account.\\
\hline
\end{tabular}
\caption{Analysis of the FTC MetaMask phishing example}
\label{tab:metamask-analysis}
\end{table}

\textbf{Why it should not be trusted:} The message combines impersonation,
urgency, and a request to click a link. A safer approach is to open the
official service independently or use a known contact method instead of
using information supplied by the message.

\referenceblock{FTC, ``Those urgent emails from MetaMask and PayPal are phishing scams.''}{https://consumer.ftc.gov/comment/181991}

\subsection{Public Example 2: Fake Direct-Deposit Phishing Email}

NIST's phishing-awareness material includes a sample direct-deposit
phishing email. The example appears to come from an individual and
claims that the sender has changed banks and needs direct-deposit
information updated before the previous account is closed. It creates
pressure to provide assistance promptly.

\begin{table}[H]
\centering
\small
\renewcommand{\arraystretch}{1.35}
\begin{tabular}{|p{4.0cm}|p{9.2cm}|}
\hline
\textbf{Indicator} & \textbf{Analysis}\\
\hline
Suspicious sender & The educational example shows an email from an address that does not provide sufficient assurance that the sender is really the person represented. Sender identity should be independently verified.\\
\hline
Suspicious link & The published description does not identify a specific embedded URL as the main indicator. The absence of a visible link does not make the request trustworthy.\\
\hline
Urgent / threatening language & The sender says the previous bank account will be closed in a few days and asks for prompt assistance, creating time pressure.\\
\hline
Spelling / grammar & The example is relatively well written. This demonstrates that correct spelling and grammar do not prove that an email is legitimate.\\
\hline
Request for personal information & The message attempts to obtain information related to direct-deposit banking details, which is highly sensitive.\\
\hline
Other suspicious element & A request to change payment or banking information should be independently verified through a trusted organizational process.\\
\hline
\end{tabular}
\caption{Analysis of the NIST direct-deposit phishing example}
\label{tab:direct-deposit-analysis}
\end{table}

\textbf{Why it should not be trusted:} A request involving salary or
banking information is high impact. The recipient should verify the
request using a known phone number, internal system, or previously
established contact method rather than replying to the suspicious
message.

\referenceblock{NIST, ``Phishing Cues-SOUPS Poster'' -- sample direct-deposit phishing email.}{https://tsapps.nist.gov/publication/get_pdf.cfm?pub_id=956178}

\subsection{Comparison of the Two Phishing Examples}

\begin{table}[H]
\centering
\small
\renewcommand{\arraystretch}{1.35}
\begin{tabular}{|p{3.0cm}|p{5.0cm}|p{5.0cm}|}
\hline
\textbf{Indicator} & \textbf{Example 1: MetaMask} & \textbf{Example 2: Direct Deposit}\\
\hline
Impersonation & Trusted cryptocurrency service & Individual / organizational identity\\
\hline
Pressure & Act quickly to avoid loss of cryptocurrency & Act promptly before an account is closed\\
\hline
Sensitive target & Wallet/account information & Banking and payroll information\\
\hline
Link involvement & Link requested to resolve the fake problem & No specific embedded link identified in the published description\\
\hline
Safe response & Independently verify the wallet/service & Independently verify through a trusted organizational channel\\
\hline
\end{tabular}
\caption{Comparison of two publicly documented phishing examples}
\label{tab:phishing-comparison}
\end{table}

\subsection{Cybersecurity Connection}

Phishing is directly connected to identity theft, credential theft,
malware delivery, financial fraud, and unauthorized account access.
A successful phishing attempt can also become the first stage of a
larger attack. Security awareness therefore acts as an important
preventive control alongside technical controls such as spam filters,
endpoint security, multi-factor authentication, and email protection.

\subsection{Security Risks and Basic Protection}

\textbf{Major Risks:}
\begin{itemize}
    \item Theft of usernames, passwords, OTPs, and financial information.
    \item Installation of malware through malicious links or attachments.
    \item Account takeover after stolen credentials are reused.
    \item Financial loss and identity theft.
    \item Further compromise of an organization after an employee is deceived.
\end{itemize}

\textbf{Basic Protection:}
\begin{itemize}
    \item Do not click unexpected links or open unexpected attachments.
    \item Verify suspicious requests using an independently known contact method.
    \item Use multi-factor authentication on important accounts.
    \item Report suspicious messages through the appropriate reporting mechanism.
    \item Delete phishing messages after reporting and documenting them when required.
\end{itemize}

\referenceblock{CISA, ``Avoid Phishing Scams with Three Simple Tips.''}{https://www.cisa.gov/sites/default/files/2024-09/Secure-Our-World-Phishing-Tip-Sheet.pdf}

\paperblock{NIST, ``Phishing'' -- Small Business Cybersecurity Guidance}{https://www.nist.gov/itl/smallbusinesscyber/guidance-topic/phishing}

\videoblock{1}{Professor Messer -- Phishing and Social Engineering}{https://www.youtube.com/results?search_query=Professor+Messer+phishing+social+engineering}

\subsection{Quick Check}

\begin{enumerate}
    \item Why should a suspicious link not be opened during analysis?
    \item Name two common warning signs of a phishing message.
    \item Why should a request for banking or password information be verified independently?
    \item Is poor grammar required for a message to be phishing?
\end{enumerate}

\subsubsection*{Answers}

\begin{enumerate}
    \item Because the link may lead to a fraudulent website or malicious content.
    \item Examples include an unusual sender, suspicious link, urgent language, unexpected attachment, or request for sensitive information.
    \item Because the message may be impersonating a trusted person or organization and could be attempting to steal sensitive information.
    \item No. A phishing message can be professionally written, so grammar alone is not a reliable indicator.
\end{enumerate}

\subsection{Task 3 Summary}

The two public examples demonstrate that phishing messages often rely
on impersonation, urgency, unexpected requests, suspicious links, and
attempts to obtain sensitive information. A safe analysis does not
require interacting with the suspicious content. Instead, the message
should be examined for warning signs and verified through an
independent trusted channel.

% ============================================================
% TASK 4
% ============================================================

\section{Task 4: Malware Prevention}

\subsection{Introduction}

Malware prevention combines technical controls, safe software habits,
secure account practices, backups, and user awareness. No single
control eliminates every malware risk, so several protective measures
should be used together.

\subsection{Malware Prevention -- Overall Security Process}

A basic prevention process begins with keeping the device updated and
using security software. Trusted applications, safe downloads,
careful handling of attachments, strong authentication, and backups
further reduce the likelihood or impact of malware incidents.

\begin{figure}[H]
\centering
\includegraphics[width=0.36\textwidth]{images/malware_prevention_overall.png}
\caption{Malware Prevention -- Overall Security Process}
\label{fig:malware-overall}
\end{figure}
\FloatBarrier

\subsection{Ten Malware Prevention Practices}

The following ten practices provide a practical baseline for reducing
malware risk on computers and mobile devices.

\begin{figure}[H]
\centering
\includegraphics[width=0.70\textwidth]{images/ten_malware_prevention_practices.png}
\caption{Ten Malware Prevention Practices}
\label{fig:ten-malware-practices}
\end{figure}
\FloatBarrier

\begin{enumerate}
    \item \textbf{Keep software updated} -- Install operating-system and application security updates.
    \item \textbf{Use security software} -- Use reputable antivirus or endpoint protection where appropriate.
    \item \textbf{Use trusted applications} -- Install software from official or otherwise trusted sources.
    \item \textbf{Use safe downloads} -- Avoid unknown files and suspicious download websites.
    \item \textbf{Be careful with attachments} -- Do not open unexpected or suspicious files.
    \item \textbf{Use strong passwords} -- Use unique passwords or passphrases for important accounts.
    \item \textbf{Enable multi-factor authentication} -- Add an additional verification factor where supported.
    \item \textbf{Maintain regular backups} -- Keep protected copies of important data.
    \item \textbf{Use firewall protection} -- Control unnecessary network connections and exposure.
    \item \textbf{Maintain security awareness} -- Recognize suspicious files, links, prompts, and behavior.
\end{enumerate}

\subsection{Detailed Explanation of Five Practices}

\subsubsection{1. Keep Software Updated}

Software updates frequently contain security fixes that address known
vulnerabilities. Keeping the operating system, browser, applications,
and security tools updated reduces exposure to vulnerabilities that
attackers may attempt to exploit.

\textbf{Practical action:} Enable automatic updates where appropriate
and regularly check for updates for software that does not update
automatically.

\subsubsection{2. Use Security Software}

Reputable antivirus or endpoint-security software can help detect,
block, quarantine, or remove malicious software. Security software is
not a replacement for safe behavior, but it provides an additional
layer of protection.

\textbf{Practical action:} Keep security software active and updated,
and perform scans when suspicious behavior is observed.

\subsubsection{3. Use Trusted Applications and Safe Downloads}

Applications should be obtained from official stores, known vendors,
or other trusted sources. Unofficial or modified installers can contain
malware or unwanted software. Users should also verify the publisher
and download source before installation.

\textbf{Practical action:} Check the official website or trusted app
store, verify the publisher, and avoid suspicious download portals.

\subsubsection{4. Avoid Unknown Attachments}

Malware can be delivered through malicious email attachments and other
unexpected files. A file that appears to be an invoice, document, or
image may contain malicious content or direct the user to a harmful
website.

\textbf{Practical action:} Verify unexpected attachments through a
separate trusted channel before opening them.

\subsubsection{5. Maintain Regular Backups}

Backups reduce the impact of data loss caused by malware, accidental
deletion, hardware failure, or ransomware. A useful backup strategy
keeps copies protected from unauthorized modification or deletion.

\textbf{Practical action:} Back up important files regularly and
periodically verify that the backup can actually be restored.

\subsection{Malware Backup and Recovery Process}

Backups are particularly important when malware causes data loss or
makes files unavailable. A safe recovery process isolates the affected
device, removes or rebuilds the compromised system, and restores data
from a trusted backup.

\begin{figure}[H]
\centering
\includegraphics[width=0.32\textwidth]{images/malware_backup_recovery.png}
\caption{Malware Backup and Recovery Process}
\label{fig:backup-recovery}
\end{figure}
\FloatBarrier

\textbf{Important point:} A backup that remains permanently connected
and writable by the infected system may also be affected by malware.
Backups should therefore be protected against unauthorized modification
or deletion.

\subsection{Safe Software Download and Installation}

Software should be obtained from a trusted source and checked before
installation. If the source is suspicious or cannot be verified, the
safest action is to stop rather than continue with the installation.

\begin{figure}[H]
\centering
\includegraphics[width=0.48\textwidth]{images/safe_software_download.png}
\caption{Safe Software Download and Installation}
\label{fig:safe-download}
\end{figure}
\FloatBarrier

\referenceblock{CISA, ``Secure Our World -- Malware and software security guidance.''}{https://www.cisa.gov/secure-our-world}

\paperblock{NIST, ``Cybersecurity Framework (CSF) 2.0'' -- protection, detection, response, and recovery concepts}{https://www.nist.gov/cyberframework}

\videoblock{2}{Professor Messer -- Malware and Malware Prevention}{https://www.youtube.com/results?search_query=Professor+Messer+malware+prevention}

\subsection{Cybersecurity Connection}

Malware prevention is a layered security activity. Updates reduce
known vulnerabilities, security software provides detection, trusted
software sources reduce supply-chain and malicious-download risks,
backups support recovery, firewalls reduce unnecessary exposure, and
security awareness reduces unsafe user actions.

These controls correspond to the broader cybersecurity cycle of
protecting systems, detecting suspicious activity, responding to
incidents, and recovering from security events.

\subsection{Security Risks and Basic Protection}

\begin{table}[H]
\centering
\small
\renewcommand{\arraystretch}{1.35}
\begin{tabular}{|p{4.0cm}|p{4.5cm}|p{4.5cm}|}
\hline
\textbf{Risk} & \textbf{Possible Impact} & \textbf{Basic Protection}\\
\hline
Unpatched software & Exploitation of known vulnerabilities & Apply security updates\\
\hline
Malicious download & Malware infection & Use trusted sources and verify software\\
\hline
Malicious attachment & Malware execution or credential theft & Verify unexpected attachments before opening\\
\hline
Weak authentication & Account compromise & Use strong unique passwords and MFA\\
\hline
Data loss / ransomware & Loss or encryption of important files & Maintain protected backups\\
\hline
Unnecessary network exposure & Increased attack surface & Use firewall controls and secure configurations\\
\hline
User error & Unsafe execution or disclosure & Maintain security awareness and training\\
\hline
\end{tabular}
\caption{Common malware risks and basic protection}
\label{tab:malware-risks}
\end{table}

\subsection{Mini Practical: Safe Software Verification}

A simple practical activity can be performed without downloading
unknown software.

\begin{enumerate}
    \item Select a commonly used application that is already installed or that you legitimately need.
    \item Locate the official vendor website or trusted application store.
    \item Check the publisher or developer name.
    \item Compare the software name and version information with the official source.
    \item Do not install the application if the source, publisher, or download details appear suspicious.
\end{enumerate}

\textbf{Security Note:} The practical should not involve executing an
unknown file or visiting a suspicious download site. The objective is
to practice safe verification, not to test malware.

\referenceblock{CISA, ``Avoiding Social Engineering and Phishing'' -- safe handling of suspicious content.}{https://www.cisa.gov/topics/cyber-threats-and-advisories/social-engineering}

\subsection{Quick Check}

\begin{enumerate}
    \item Why are software updates important for malware prevention?
    \item Why should software be downloaded from trusted sources?
    \item How can backups reduce the impact of ransomware or data loss?
    \item Why are strong passwords and MFA useful against malware-related account compromise?
    \item Why should unexpected attachments be verified before opening?
\end{enumerate}

\subsubsection*{Answers}

\begin{enumerate}
    \item Updates can fix known security vulnerabilities that attackers may exploit.
    \item Trusted sources reduce the chance of installing modified, malicious, or unwanted software.
    \item Backups provide clean copies of important data that can be restored after an incident.
    \item They make unauthorized account access more difficult, even if a password is exposed.
    \item An unexpected attachment may contain malicious content or lead to further compromise.
\end{enumerate}

\subsection{Task 4 Summary}

Malware prevention requires multiple layers of protection rather than
one security tool. Keeping software updated, using security software,
choosing trusted applications and downloads, handling attachments
carefully, protecting accounts, maintaining backups, using firewalls,
and developing security awareness all contribute to reducing malware
risk.

% ============================================================
% DOCUMENTATION REFERENCES
% ============================================================

\Needspace{5cm}
\subsection{Documentation References}

\begin{enumerate}
    \item CISA -- Secure Our World and phishing-awareness guidance.
    \item FTC Consumer Advice -- phishing and scam-awareness guidance.
    \item NIST -- Phishing guidance and phishing-awareness examples.
    \item NIST Cybersecurity Framework 2.0.
\end{enumerate}

\Needspace{4cm}
\subsection{YouTube Learning References}

\begin{enumerate}
    \item Professor Messer -- Phishing and Social Engineering.
    \item Professor Messer -- Malware and Malware Prevention.
\end{enumerate}

\end{document}
